\chapter{ML Systems and Production}
\label{chap:systems}
A trained model file is one component inside a data, serving and feedback loop.
Production quality therefore depends on data freshness, feature consistency, serving reliability, monitoring and retraining as well as the predictor itself.

\section{The full production loop}
\label{sec:systems-production-loop}
The useful system picture is
\[
\boxed{
\text{events}
\rightarrow
\text{features}
\rightarrow
\text{prediction}
\rightarrow
\text{outcome/logs}
\rightarrow
\text{future training data}
}.
\]

\begin{figure}[htbp]
\centering
\resizebox{.97\linewidth}{!}{%
\begin{tikzpicture}[node distance=8mm and 9mm]
\node[flow] (events) {World / user\\events};
\node[flow,right=of events] (ingest) {Ingestion};
\node[flow,right=of ingest] (online) {Online\\features};
\node[flow,right=of online] (serve) {Model\\serving};
\node[yesbox,right=of serve] (pred) {Prediction};
\node[flow,below=14mm of ingest] (store) {Historical\\store};
\node[flow,right=of store] (train) {Training\\dataset};
\node[flow,right=of train] (fit) {Model\\training};
\node[yesbox,right=of fit] (registry) {Validation +\\registry};
\draw[arr] (events)--(ingest);\draw[arr] (ingest)--(online);\draw[arr] (online)--(serve);\draw[arr] (serve)--(pred);
\draw[arr] (ingest)--(store);\draw[arr] (store)--(train);\draw[arr] (train)--(fit);\draw[arr] (fit)--(registry);
\draw[arr,rounded corners] (registry.east)--++(.8,0)--([xshift=.8cm]serve.south)--(serve.south);
\draw[arr,rounded corners] (pred.south)--++(0,-1.7)--([yshift=-1.7cm]store.east)--(store.east);
\end{tikzpicture}}
\caption{A production ML system closes the loop from observations to predictions and back into future training data.}
\end{figure}

\section{Batch versus streaming}
\label{sec:systems-batch-streaming}
Batch processing waits for a collection of events and processes them together.
Streaming updates state as events arrive.
The core trade-off is
\[
\underbrace{\text{freshness}}_{\text{streaming advantage}}
\quad\longleftrightarrow\quad
\underbrace{\text{operational simplicity}}_{\text{batch advantage}}.
\]
If a feature changes slowly relative to the decision cadence, streaming adds complexity without much value.
If session intent becomes useful within seconds, waiting for tomorrow's batch can make the feature operationally irrelevant.

\begin{calloutpractical}
Ask how quickly new information changes the optimal decision before deciding that a feature needs a streaming path.
\end{calloutpractical}

\section{Training-serving consistency}
\label{sec:systems-training-serving-skew}
A model is validated on training-time feature semantics.
If serving computes the ``same'' named feature differently, deployment receives a distribution the model was never validated on.
A feature store or shared transformation layer mainly helps by keeping
\[
\boxed{
\text{training definition}
\approx
\text{serving definition}
}
\]
for code, windows, units and missing-value behaviour.
The same leakage discipline from \cref{sec:evaluation-model-selection-and-data-leakage} applies to historical feature construction.

\section{Model registry, deployment and rollback}
\label{sec:systems-registry}
A deployed prediction should be traceable to the model version, training data/version, feature schema, code/configuration and evaluation report that produced it.
A registry makes that state explicit and supports rollback when the new version fails after deployment.
The practical unit is therefore not ``the latest weights'' but a versioned model artifact plus its reproducible context.

\section{Monitoring as a stack}
\label{sec:systems-monitoring}
One ``model health'' metric cannot diagnose infrastructure, data and product failures at once.

\begin{figure}[htbp]
\centering
\begin{tikzpicture}
\draw[deepblue,thick,fill=palerblue] (0,0) rectangle (10,1.2);
\node at (5,.6) {\textbf{Infrastructure:} latency, errors, accelerator memory};
\draw[teal,thick,fill=palerteal] (0,1.35) rectangle (10,2.55);
\node at (5,1.95) {\textbf{Data:} schema, missingness, feature distributions};
\draw[amber,thick,fill=paleamber] (0,2.7) rectangle (10,3.9);
\node at (5,3.3) {\textbf{Model:} score distribution, calibration, slice metrics};
\draw[navy,thick,fill=white] (0,4.05) rectangle (10,5.25);
\node at (5,4.65) {\textbf{Product:} clicks, revenue, safety outcomes, retention};
\end{tikzpicture}
\caption{Monitoring layers answer different questions; a healthy server can host a failing model and a good model can still time out.}
\end{figure}

\section{Drift and retraining}
\label{sec:systems-drift-retraining}
Covariate drift changes $P(x)$; concept drift changes $P(y\mid x)$.
A moved feature histogram does not prove model quality fell, and model quality can fall without an obvious one-dimensional drift signal.
Retraining should therefore follow diagnosis rather than act as an automatic cure for every drift alert.

Common policies are:
\[
\text{fixed schedule}
\qquad
\text{event-triggered retraining}
\qquad
\text{frequent incremental updates}.
\]
More frequent updates improve freshness but increase reproducibility, stability and feedback-loop complexity.

\section{Queues, caching and fallbacks}
\label{sec:systems-fallbacks}
Production systems must specify what happens when the preferred model is slow or unavailable.
A fallback can be a cached result, a smaller model, a heuristic ranking or a safe default.
Slightly weaker predictions can be operationally preferable to request failure.

\begin{figure}[htbp]
\centering
\begin{tikzpicture}[node distance=8mm and 10mm]
\node[flow] (req) {Request};
\node[flow,right=of req] (main) {Primary model};
\node[decisionnode,right=of main] (ok) {success\\on time?};
\node[yesbox,above right=5mm and 12mm of ok] (out) {Return result};
\node[flow,below right=5mm and 12mm of ok] (fall) {Fallback\\model/cache/rule};
\draw[arr] (req)--(main);\draw[arr] (main)--(ok);
\draw[arr] (ok)--node[above,font=\scriptsize]{yes}(out);
\draw[arr] (ok)--node[below,font=\scriptsize]{no}(fall);
\end{tikzpicture}
\caption{Failure handling is part of the prediction policy, not an afterthought.}
\end{figure}

\section{Build versus buy}
\label{sec:systems-build-buy}
API serving and self-hosting have different fixed/marginal cost structures:
\[
\underbrace{C_{\text{lifetime}}}_{\text{total}}
=
\underbrace{C_{\text{fixed engineering + infrastructure}}}_{\text{larger for self-hosting}}
+
\underbrace{N_{\text{requests}}C_{\text{request}}}_{\text{often larger for API use}}.
\]
The break-even volume is not universal because utilisation, token lengths, hardware, staffing, privacy and model choice change both terms.

\begin{center}
\small
\begin{tabularx}{.94\linewidth}{@{}L{.27\linewidth}YY@{}}
\toprule
\textbf{Dimension} & \textbf{API} & \textbf{Self-host} \\
\midrule
Launch speed & Usually high & Infrastructure first \\
Operations burden & Provider owns more & Team owns serving stack \\
Up-front cost & Lower & Higher \\
Model/serving control & Limited by provider & High \\
Marginal economics & Provider pricing & Depends on utilisation/hardware \\
Privacy/control & Provider-dependent & Stronger direct control possible \\
\bottomrule
\end{tabularx}
\end{center}

\section{Training versus lifetime inference economics}
\label{sec:systems-inference-economics}
Training is paid a small number of times; inference can be paid for every user request and every generated token.
Lifecycle cost is therefore
\[
\boxed{
C_{\text{lifetime}}
=
C_{\text{training}}
+
N_{\text{inferences}}
C_{\text{inference}}.
}
\]
Extra training compute can be rational if it produces a smaller model that reaches the required quality and is then served at enormous volume.
The interview-guide heuristic cites Chinchilla-style compute-optimal training at roughly 20 training tokens per parameter as a historical reference point.
Treat that ratio as a result for a particular modelling/compute regime rather than a universal modern constant.
Do not turn a historical scaling-law ratio into a universal constant; the useful principle is the one-off training versus repeated-serving trade-off.

\begin{calloutmemoryhook}
Training is paid occasionally; serving is paid every request.
\end{calloutmemoryhook}

\begin{calloutdegreeofexperiment}
Vary batch versus streaming paths, feature freshness, serving model, caching, fallback policy, deployment cadence, retraining trigger, observability thresholds and build-versus-buy boundary.
\end{calloutdegreeofexperiment}
